\section{Bounded full Brauer identities and explicit contraction diagrams}
\label{sec:e8-transfer-diagrams}

\paragraph{Scope and provenance.}
This section proves the representation-theoretic and correspondence-count
steps needed for the quantitative branch-exponent construction.  The
integer column-lattice argument is reconstructed from the proof inside
the supplied SQ3.  The printed statement of SQ3 concerns determinant
characters; it is \emph{not} silently used as a bounded identity of
representations.  The latter identity is proved below.  The geometric
supplier theorem HD-FS1--02 and the separately bounded Fermat-product
Hodge cycles remain explicit inputs.  No arbitrary endomorphism of a
multiplicity Hodge structure is asserted to be algebraic.

Let \(\mathcal G\) be the connected-cover group, write \(G=|\mathcal G|\),
and let \(C\) be its smooth projective curve of genus \(g\geq1\).
Let \(F\subset\mathbb C\) be a finite Galois splitting field for
\(\mathcal G\) containing all \(G\)-th roots of unity, and put
\(A=F[\mathcal G]\).  This ensures that every subgroup's linear
characters used below are defined over \(F\); splitting
\(\mathcal G\) alone would not suffice for that assertion.
The coefficient representation \(V\) is absolutely irreducible.
In the application \(G=1393459200\), \(g=1+53G/24=3077222401\),
and the multiplicity space of \(V\) in \(H^1(C,F)\) has dimension six.

\subsection*{D-E8-XFER-01. A bounded identity of honest representations}

There is an integral Brauer identity
\[
 [V]=\sum_j n_j[\operatorname{Ind}_{H_j}^{\mathcal G}\chi_j],
 \qquad
 \sum_j|n_j|\leq (G+1)G!G^G
 \leq R:=(2G+1)G!G^G.                                  \tag{XD.1}
\]
Here each \(\chi_j\) is a linear character and trivial characters are
allowed.  Thus, on moving negative terms to the other side, (XD.1)
is an isomorphism
\[
 V\oplus R_-\simeq R_+                                   \tag{XD.2}
\]
of actual \(F[\mathcal G]\)-modules, with at most \(R\) induced
summands counted with their repetitions.

\begin{proof}
Use all irreducible representation multiplicities as rows of an
integer matrix.  For every pair \((H,\chi)\), take the multiplicity
vector of \(\operatorname{Ind}_H^{\mathcal G}\chi\) as a column.
There are finitely many columns, at most \(G\) rows, and every
entry lies between zero and \(G\): it is at most the dimension of
the induced representation, which is \([\mathcal G:H]\leq G\).
Integral Brauer induction puts the unit vector for \(V\) in the
integral column lattice.

For clarity, here is the needed integer bound.  Suppose an integer
column matrix has rank \(r\), entries bounded by \(M\geq1\), and a
target \(b\), with \(\|b\|_\infty\leq1\), in its integral column
lattice.  Choose \(r\) independent rows and \(r\) independent
columns forming an invertible square matrix \(B\).  Projection to
these rows is injective on the rational column span.  Set
\[
                    d=|\det B|\leq r!M^r.
\]
The subgroup generated by the remaining columns in
\(\mathbb Z^r/B\mathbb Z^r\) has at most \(d\) elements.  A shortest
path from zero to the class of \(b\), using these columns and
their negatives, has at most \(d-1\) edges.  Subtract the corresponding
signed columns.  The residual vector \(v\in B\mathbb Z^r\) satisfies
\[
                  \|v\|_\infty\leq1+(d-1)M\leq dM.
\]
Every cofactor is bounded by \((r-1)!M^{r-1}\).  Cramer's rule,
with its denominator \(d\), therefore gives
\[
              |(B^{-1}v)_j|\leq r!M^r.
\]
The total absolute coefficient sum is at most
\((r+1)r!M^r\).  The identity in the selected rows is an identity
in all rows by injectivity on the column span.  Taking \(r\leq G\)
and \(M=G\) proves (XD.1).  Semisimplicity in characteristic zero
turns equality of representation classes into (XD.2).
\end{proof}

Write \(U=H^1(C,F)\), let \(M_\pm=\operatorname{Hom}_{\mathcal G}(R_\pm,U)\),
and suppose \(\dim_F\operatorname{Hom}_{\mathcal G}(V,U)=6\).
Put \(k=\dim_F M_-\).  Each induced summand contributes a subspace
of \(U\), of dimension at most \(2g\).  Applying the contravariant
multiplicity functor to (XD.2) gives
\[
 M_+\simeq M_V\oplus M_-,\qquad
 \ell:=\dim_F M_++\dim_F M_-=6+2k\leq2gR=:L.                \tag{XD.3}
\]
Zero-dimensional summands can be omitted \emph{after} applying this
functor; no claim is made that deleting them preserves (XD.2) as a
representation identity.  In particular the actual curve-slot count
\(\ell\) is even and at least six.

\subsection*{D-E8-XFER-02. Every required map is a finite sum of deck diagrams}

Fix a very ample divisor \(H_C\) of degree \(6g\).  The homomorphism
of determinant lines obtained from (XD.3) and the polarization,
\[
             \det M_+\otimes\det\overline{M_-}
                  \longrightarrow\det M_V                 \tag{XD.4}
\]
is nonzero.  After embedding all its factors in tensor powers of
\(U\), it is a finite \(F\)-linear combination of the following maps:
apply deck transformations to individual inputs, pair \(k\) disjoint
pairs of inputs by the cup-product pairing, and send the six remaining
inputs, in some order, to six output factors.  Every such summand is
realized by an integral effective correspondence
\[
                     Z\subset C^\ell\times C^6
\]
which is, after permuting coordinates, a product of \(k+6\) deck
graphs.  Its degree for the sum of the curve polarizations is
\[
                 \deg Z=(k+6)!(12g)^{k+6}.                  \tag{XD.5}
\]
Precomposing with the faithful cyclic-character and determinant
projectors introduces only further sums of the same kind.

\begin{proof}
This is an assertion about maps induced by representations, not
about all Hodge maps on \(U\).  For a linear character define
\[
 e_{H,\chi}=\frac1{|H|}\sum_{h\in H}\chi(h)^{-1}h\in A.
\]
The left ideal \(Ae_{H,\chi}\) is the induced module
\(\operatorname{Ind}_H^{\mathcal G}\chi\), and
\[
 \operatorname{Hom}_A(Ae,U)\simeq eU,\qquad
                     f\longmapsto f(e).                    \tag{XD.6}
\]
Choose a primitive idempotent \(e_V\) with \(Ae_V\simeq V\).
In the \(E_8\) application choose
\(e_V=\epsilon_V(1-s_1)/2\), where \(\epsilon_V\) is the
individual central idempotent for the chosen embedding of the
natural character.  Then \(e_V+\overline{e_V}\) is precisely
the rational root-line projector \(p\) used in the last map.
If another primitive idempotent is first used, its fixed
module isomorphism with this one is also an entry of the
group algebra by the next argument.
The isomorphism (XD.2) is therefore an isomorphism between direct
sums of left ideals of the semisimple algebra \(A\).

Every \(A\)-linear map \(Ae\to Af\) is determined by the image of
\(e\), an element \(t\in eAf\); explicitly it sends \(ae\) to \(at\).
After applying (XD.6), its contravariant action is \(u\mapsto tu\).
Consequently the induced maps in (XD.3) are matrices whose entries
are elements of \(A\).  Each entry is a finite \(F\)-linear
combination of actual deck transformations of \(C\), and is
realized by the corresponding combination of their graphs.
This gives the required algebraic maps without a density assertion
about unknown multiplicity-Hodge endomorphisms.

The cup product is a nondegenerate alternating form \(Q\) on \(U\).
The transpose of \(e_{H,\chi}\) for this form is
\(e_{H,\chi^{-1}}\).  Thus it pairs \(e_{H,\chi}U\)
nondegenerately with \(e_{H,\chi^{-1}}U\), and the direct-sum
pairing identifies \(\overline{M_-}\) with \(M_-^*\), up to the
usual Tate twist.  Take the top exterior power of (XD.3) and then
contract these two determinant factors.  The isomorphism and the
nondegenerate determinant pairing prove that (XD.4) is nonzero.

More explicitly, write the two components of
\(M_+\to M_V\oplus M_-\) as \(a\) and \(b\).  For each choice of
six of the \(6+k\) plus inputs, send those six through \(a\);
send the others through \(b\) and pair them with the \(k\) minus
inputs.  Alternate over the choices, the six outputs and the
minus-input matching.  This is the usual exterior-power
isomorphism followed by the determinant pairing.  Expanding
each matrix entry in \(F[\mathcal G]\) gives precisely the
asserted maps.  The factors from direct-sum labels are just
scalar coefficients and choices of the corresponding entries.

For paired slots, group invariance gives
\(Q(gx,hy)=Q(x,g^{-1}hy)\); hence one deck graph on the two
input curves realizes each pairing.  A retained input and its
output use one deck graph as well.  The \(\ell+6=2(k+6)\)
coordinates are thereby partitioned into disjoint pairs.
The graph on each pair is isomorphic to \(C\) and has degree
\(2\deg H_C=12g\).  The product polarization gives (XD.5).
This also verifies the correspondence's degree shift:
its codimension is \((\ell+6)/2\), so it sends \(H^\ell(C^\ell)\)
to \(H^6(C^6)\).

The cyclic-character projectors act coordinatewise by group-algebra
elements.  Determinant alternation acts by permutations, with
the Koszul sign when it is written in cohomological tensors.
Precomposition by either operation merely changes the deck words,
the input pairing, the output ordering, or their scalar coefficients.
It remains a sum of the same diagrams.
\end{proof}

There is one further elementary projector needed when cycles have
components in other Kunneth degrees.  For a chosen point \(o\in C\),
\[
 \pi_1=[\Delta_C]-[\{o\}\times C]-[C\times\{o\}]
\]
projects onto \(H^1(C)\).  Every term of \(\pi_1^{\boxtimes\ell}\)
is a product of \(\ell\) diagonals or point-times-curve factors.
It is an integral effective cycle of dimension \(\ell\) and degree
at most
\[
                         \ell!(12g)^\ell.                   \tag{XD.7}
\]
Thus the full tensor Kunneth projection is one correspondence
operation, expanded into finitely many terms.  The number or sizes
of their scalar coefficients are not a degree cost for the
existence conclusion below.

\subsection*{D-E8-XFER-03. Trivial characters and bounded component selection}

A trivial inducing character contributes \(H^1(B,F)\) for
\(B=C/H\), with \(b=g(B)\).  Its determinant is represented,
after Kunneth projection and alternation, by a nonzero scalar
multiple of the product of \(b\) pullback diagonals
\[
                  T_H=(C\times_B C)^b\subset C^{2b}.
\]
As fundamental cycles \(C\times_B C=\sum_{h\in H}\Gamma_h\):
the equality holds over the free locus and both sides have
pure dimension one.  Thus
\[
                   \deg T_H=b!(12g|H|)^b
                              \leq b!(12gG)^b.              \tag{XD.8}
\]
Nonvanishing follows from the nondegenerate symplectic pairing
on \(H^1(B)\): its inverse-pairing tensor, repeated \(b\) times
and alternated, is its top symplectic volume.  Quotient pullback
only multiplies the pairing by the nonzero integer \(|H|\).

These trivial factors must be external factors of the returning
correspondence.  They are \emph{not} described as finite graphs over
a point.  If \(\ell_{\rm non}\) slots belong to the nontrivial
characters and \(\ell_{\rm triv}\) to trivial characters, the
combined correspondence is
\[
                 \Gamma_{\rm non}\times T_{\rm triv}
                         \subset Y\times C^\ell.            \tag{XD.9}
\]
If \(d=\dim Y\), its dimension is \(d+\ell_{\rm triv}/2\).
It sends a source class of degree \(2d-\ell_{\rm non}\) to a
class of degree \(\ell\), as required.  The factors with \(b=0\)
are units and are omitted.

For clarity, the following statement isolates the quantitative
selection step.  Suppose every ambient product used here is smooth
projective of dimension at most \(2L\), with very ample product
polarization of degree at most \(A\), and every required returning
correspondence (XD.9) has degree at most \(A\).
Suppose the relevant rational Hodge classes on each \(Y\) are spanned
by integral effective cycles of degree at most \(S\).
Take
\[
 A=(4L)!\,[1000g(G+1)(L+1)]^{4L},\qquad
                         B=(4L+4)A^{2L+4}.                 \tag{XD.10}
\]
The elementary degree estimates (XD.5), (XD.7) and (XD.8) are all
less than \(A\), since \(6\leq\ell\leq L\) and
\(k+6=(\ell+6)/2\leq\ell\).

The returning operation, the full tensor Kunneth projector, the
diagram operation of D-E8-XFER-02, and a final coordinatewise
primitive projector constitute at most four correspondence
operations.  All cyclic-character projectors, exterior-power
intertwiners and determinant alternations have already been
expanded into the third operation.  The final primitive projector
is redundant if included there, but keeping it costs at most one
additional operation.

Here is the degree cost for each operation, including improper
intersections.  If \(Z\) has degree \(E\), taking its product with
the target has degree at most \(AE\): the relevant binomial
coefficient times the target degree is bounded by the ambient
product degree.  Intersecting it with a correspondence of degree
at most \(A\), the effective positive-minus-negative version of
Chow moving in Truong's Lemma 2.2 has total mass at most
\[
        2(2L+1)A^{2L+1}(AE)A
             =2(2L+1)A^{2L+3}E\leq BE.                     \tag{XD.11}
\]
The displayed factor is deliberately larger than the published
constant \(kD^k\) for each sign.  Proper projection does not
increase degree for the sum polarization: the omitted source
hyperplane class is nef.  The same mass estimate therefore
holds after pushforward.

The nonzero determinant map (XD.4), the supplier theorem and the
nonzero pairing in its returning correspondence show that the
composite supplies the individual \(F\)-determinant line.
To use rational Hodge cycles, first expand all group-algebra and
permutation coefficients and let \(I\) be the sum of the images
of the resulting \emph{rational} correspondence maps from the
supplier cohomologies, with the required Tate shifts.
This is a rational Hodge substructure.  Its scalar extension
\(I_F\) contains the determinant line, hence also all its
\(\operatorname{Gal}(F/\mathbb Q)\)-conjugates, because \(I\)
is rational.  The conjugate pair spans the rational Weil
plane after scalar extension, so \(I\) contains that Weil
plane over \(\mathbb Q\).
Semisimplicity of polarizable rational Hodge structures now
gives a \emph{rational Hodge preimage} of a nonzero rational
Weil class in the direct sum of the supplier cohomologies.
The assumed rational spanning cycles express this preimage
as a finite linear combination.  Linearity implies that
some single initial spanning cycle and some choices of diagram
terms have nonzero final Weil projection.  At each moving
intersection write the result as integral positive minus negative
cycles using (XD.11).  Expanding these finite sums again, at least
one individual integral effective component has nonzero final
Weil projection.  Here the projection may be evaluated after the
fixed last map to \(B\) described below; it is not asserted that
an individual term still belongs to every intermediate projector
image.  Its degree on \(C^6\) is at most
\[
                              B^4S.                         \tag{XD.12}
\]
This argument remains valid after extending the coefficient field:
a nonzero linear combination over \(F\) cannot be formed from
vectors all having zero Weil projection.  No denominator clearing,
integral saturation, or bound on a scalar coefficient is used.

Finally let \(p\) be the Rosati-fixed primitive projector and
\(M=2G\) an integer clearing it.  The homomorphism \(Mp:J(C)\to B\),
composed with Abel maps and addition, gives \(C^6\to B\).
The pullback of the conjugate Weil determinant has, on the six
curve factors, the nonzero alternating tensor of the six-dimensional
conjugate character eigenspace.  Its cup pairing with the output
determinant in (XD.4) is nonzero by the perfect determinant pairing.
Projection formula therefore proves that this last pushforward
is nonzero on the determinant output.  Component selection above
can consequently be made using the resulting Weil projection on
\(B\).  Since
\[
 (Mp)^*\theta_B\leq M^2\Theta_{J(C)},\qquad
 \deg(\Theta_{J(C)}|_C)=g,
\]
the Cauchy inequality for the sum of six maps and
\(\deg H_C=6g\) give pullback polarization at most \(M^2H_{C^6}\).
A surviving three-dimensional cycle consequently has degree at most
\[
                         M^6B^4S
                       \leq M^6B^{4L+21}S.                 \tag{XD.13}
\]
The latter is the deliberately conservative exponent used in the
main numerical theorem.  Its source-degree conclusion is conditional
only on the separately verified incidence bound \(A\) and supplier
spanning bound \(S\), not on an uncounted sequence of tensor
intertwiners.
